\chapter*{Preface}
\addcontentsline{toc}{chapter}{Preface}
\markboth{Preface}{Preface}

This book develops a cosmological framework from a single organizing principle:
\emph{physics is continuous and conservative, so singularities and
discontinuities must be artifacts of incomplete theory, not real features of
nature.} Applied carefully---together with quantum field theory, special
relativity, Einstein--Cartan gravity, and an infinite quantum vacuum---this
principle produces a complete cosmological framework with no fine-tuning. I call
it \emph{Supraverse Cartasis Theory} (\SCT). \Cartasis is the bounce
membrane: the hypersurface in spacetime where torsion-mediated repulsion
overwhelms gravitational attraction and the geometry transitions between a
parent region and a child universe.

\section*{The one-paragraph version}

Take quantum mechanics, special relativity, and Einstein--Cartan gravity (the
natural extension of general relativity that allows torsion to couple to
fermion spin). Add an infinite quantum vacuum. The vacuum is unstable: rare
large fluctuations reach a critical density (the Cartasis density,
$\rhoc \sim 10^{50}\,\mathrm{kg\,m^{-3}}$) where torsion-mediated repulsion
overwhelms gravity, and the geometry bounces rather than producing a
singularity. The bounce creates a baby universe connected to its parent region
through a 3D hypersurface, the Cartasis membrane. The supraverse---the total 4D
manifold---condenses into a foam of universes through this mechanism, with equal
statistical numbers of matter-biased and antimatter-biased lineages (CPT
symmetry preserved globally). Our universe is one bubble in this foam. The CMB
is our parent black hole's Hawking radiation, redshifted by our internal
expansion. Dark matter is the gravitational influence of our parent's matter
content (and rejected antimatter content) on us through the Cartasis membrane.
The arrow of time is the thermodynamic gradient from our bounce (low entropy)
toward our eventual Hawking-mediated dissolution back into the parent's exterior.

\section*{Why this might be worth taking seriously}

From minimal axioms the framework derives explanations for the Big Bang (a
bounce, viewed from inside), inflation (rapid torsion-driven expansion at the
bounce), the matter--antimatter asymmetry, dark matter, dark energy, the CMB,
the arrow of time, apparent fine-tuning, and information preservation---with no
singularity, no fine-tuned initial conditions, no exotic new fields, and no
special starting point. Just continuous physics applied carefully.

\section*{How the book is organized}

Chapter~\ref{chap:axioms} states the four axioms and what follows from them.
Part~II (Chapters~\ref{chap:cartasis}--\ref{chap:chirality}) develops the bounce
and its offspring. Part~III
(Chapters~\ref{chap:supraverse}--\ref{chap:timearrow}) develops the global
structure. Part~IV (Chapters~\ref{chap:tests}--\ref{chap:viz}) lays out the
tests and the computational and visualization programs.
Appendix~\ref{app:openq} catalogs what we do not yet know. Each chapter is
self-contained enough to read on its own but assumes the axioms of
Chapter~\ref{chap:axioms}. Working markdown drafts, scratch derivations, and
simulation code are kept in \texttt{drafts/}, \texttt{notes/}, and
\texttt{sims/} respectively.
